%%%% RESUMO
%%
%% Apresentação concisa dos pontos relevantes de um texto, fornecendo uma visão rápida e clara do conteúdo e das conclusões do
%% trabalho.

\begin{resumoutfpr}%% Ambiente resumoutfpr
A detecção de insetos-praga em culturas agrícolas é fundamental para o Manejo Integrado de Pragas (MIP), contudo, abordagens baseadas estritamente em visão espacial frequentemente falham diante de cenários de camuflagem crítica e alvos de pequena escala. Este Trabalho de Conclusão de Curso propõe e avalia uma arquitetura híbrida de Deep Learning baseada na rede U-Net, estendida para processar dados espácio-temporais. Para romper a ambiguidade visual imposta pelo mimetismo (onde a cor e a textura do inseto se confundem estatisticamente com o fundo foliar), o modelo incorpora a estimativa cinemática através do Fluxo Óptico Denso de Farneback, resultando em tensores de entrada de 6 canais (RGB e componentes de fluxo). Para viabilizar o treinamento sem anotação manual exaustiva, implementou-se um pipeline de supervisão fraca combinando a equalização adaptativa CLAHE e modelos de Mistura de Gaussianas (MOG2). Os resultados experimentais comprovaram que a rede U-Net espacial tradicional fracassa no cenário extremo de camuflagem verde-sobre-verde devido ao desbalanceamento e à ausência de contraste. A aplicação de aumento de dados (data augmentation) mitigou parcialmente o overfitting, porém apenas a fusão precoce das assinaturas de movimento foi capaz de segmentar consistentemente a anatomia dos espécimes em todos os cenários. Conclui-se que a incorporação do viés temporal é altamente eficaz para regularizar redes convolucionais em problemas de detecção biológica em pequena escala, oferecendo uma alternativa robusta para sistemas autônomos de diagnóstico foliar.

\singlespacing{\textcolor{red}{\footnotesize De acordo com a NBR 6028:2021, a apresentação gráfica deve seguir o padrão do documento no qual o resumo está inserido. Para definição das palavras-chave (e suas correspondentes em inglês no abstract) consultar em Termo tópico do \textbf{Catálogo de Autoridades} da Biblioteca Nacional, disponível em:} \footnotesize{http://acervo.bn.gov.br/sophia\_web/autoridade.}}
\end{resumoutfpr}

